

\olchapter{sfr}{siz}{The Size of Sets}

\begin{editorial}
This chapter discusses enumerations, countability and uncountability.
Several sections come in two versions: a more elementary one, that
takes enumerations to be lists, or surjections from $\PosInt$; and a
more abstract one that defines enumerations as bijections with $\Nat$.
\end{editorial}



\olfileid{sfr}{siz}{int}

\olsection{Introduction}

When Georg Cantor developed set theory in the 1870s, one of his aims
was to make palatable the idea of an infinite collection---an actual
infinity, as the medievals would say.  A key part of this was his
treatment of the \emph{size} of different sets. If $a$, $b$ and $c$ are
all distinct, then the set $\{a, b, c\}$ is intuitively \emph{larger}
than $\{a, b\}$. But what about infinite sets? Are they all as large
as each other? It turns out that they are not.

The first important idea here is that of an enumeration.  We can
list every finite set by listing all its !!{element}s.  For some
infinite sets, we can also list all their !!{element}s if we allow the
list itself to be infinite. Such sets are called !!{enumerable}.
Cantor's surprising result, which we will fully understand by the end
of this chapter, was that some infinite sets are not !!{enumerable}.





\olfileid{sfr}{siz}{enm}

\olsection{Enumerations and \usetoken{S}{enumerable} Sets}

\begin{editorial}
  This section discusses enumerations of sets, defining them as
  surjections from $\PosInt$. It does things slowly, for readers with
  little mathematical background. An alternative, terser
  version is given in \olref[enm-alt]{sec}, which defines enumerations
  differently: as bijections with $\Nat$ (or an initial segment).
\end{editorial}

\begin{explain}
We've already given examples of sets by listing their !!{element}s.
Let's discuss in more general terms how and when we can list the
!!{element}s of a set, even if that set is infinite.
\end{explain}

\begin{defn}[Enumeration, informally]
Informally, an \emph{enumeration} of a set~$A$ is a list (possibly
infinite) of !!{element}s of~$A$ such that every !!{element} of $A$
appears on the list at some finite position. If $A$ has an
enumeration, then $A$ is said to be \emph{!!{enumerable}}.
\end{defn}

\begin{explain}
A couple of points about enumerations:
\begin{enumerate}
\item We count as enumerations only lists which have a beginning and
  in which every !!{element} other than the first has a single
  !!{element} immediately preceding it.  In other words, there are
  only finitely many elements between the first !!{element} of the
  list and any other !!{element}. In particular, this means that every
  !!{element} of an enumeration has a finite position: the first
  !!{element} has position~$1$, the second position~$2$, etc.
\item We can have different enumerations of the same set~$A$ which
  differ by the order in which the !!{element}s appear: $4$, $1$,
  $25$, $16$,~$9$ enumerates the (set of the) first five square
  numbers just as well as $1$, $4$, $9$, $16$,~$25$ does.
\item Redundant enumerations are still enumerations: $1$, $1$, $2$,
  $2$, $3$, $3$,~\dots{} enumerates the same set as $1$, $2$,
  $3$,~\dots{} does.
\item Order and redundancy \emph{do} matter when we specify an
  enumeration: we can enumerate the positive integers beginning with
  $1$, $2$, $3$, $1$, \dots{}, but the pattern is easier to see when
  enumerated in the standard way as $1$, $2$, $3$, $4$,~\dots
\item Enumerations must have a beginning: \dots, $3$, $2$, $1$ is not
  an enumeration of the positive integers because it has no first
  !!{element}. To see how this follows from the informal definition,
  ask yourself, ``at what position in the list does the number 76
  appear?''
\item The following is not an enumeration of the positive integers:
  $1$, $3$, $5$, \dots, $2$, $4$, $6$, \dots\@ The problem is that the
  even numbers occur at places $\infty + 1$, $\infty + 2$, $\infty +
  3$, rather than at finite positions.
\item The empty set is enumerable: it is enumerated by the empty list!{}
\end{enumerate}
\end{explain}

\begin{prop}
  If $A$ has an enumeration, it has an enumeration without
  repetitions.
\end{prop}

\begin{proof}
  Suppose $A$ has an enumeration $x_1$, $x_2$, \dots{} in which each
  $x_i$ is an !!{element} of~$A$.  We can remove repetitions from an
  enumeration by removing repeated !!{element}s. For instance, we can
  turn the enumeration into a new one in which we list $x_i$ if
  it is !!a{element} of~$A$ that is not among $x_1$, \dots,
  $x_{i-1}$ or remove $x_i$ from the list if it already appears among
  $x_1$, \dots,~$x_{i-1}$.
\end{proof}

The last argument shows that in order to get a good handle on
enumerations and !!{enumerable} sets and to prove things about them,
we need a more precise definition.  The following provides it.

\begin{defn}[Enumeration, formally] 
An \emph{enumeration} of a set $A \neq \emptyset$ is any
!!{surjective} function $f \colon \PosInt \to A$.
\end{defn}

\begin{explain}
Let's convince ourselves that the formal definition and the informal
definition using a possibly infinite list are equivalent. First, any
!!{surjective} function from $\PosInt$ to a set~$A$ enumerates~$A$.
Such a function determines an enumeration as defined informally above:
the list $f(1)$, $f(2)$, $f(3)$, \dots. Since $f$ is !!{surjective},
every !!{element} of~$A$ is guaranteed to be the value of~$f(n)$ for
some~$n \in \PosInt$. Hence, every !!{element} of $A$ appears at some
finite position in the list. Since the function may not be
!!{injective}, the list may be redundant, but that is acceptable (as
noted above).

On the other hand, given a list that enumerates all !!{element}s
of~$A$, we can define !!a{surjective} function $f\colon \PosInt \to A$
by letting $f(n)$ be the $n$th !!{element} of the list, or the final
!!{element} of the list if there is no $n$th !!{element}. The only
case where this does not produce !!a{surjective} function is when $A$ is
empty, and hence the list is empty. So, every non-empty list
determines !!a{surjective} function $f\colon \PosInt \to A$.
\end{explain}

\begin{defn}
  \ollabel{defn:enumerable}
  A set~$A$ is !!{enumerable} iff it is empty or has an enumeration.
\end{defn}

\begin{ex}
A function enumerating the positive integers ($\PosInt$) is simply the
identity function given by $f(n) = n$. A function enumerating the
natural numbers $\Nat$ is the function $g(n) = n - 1$.
\end{ex}

\begin{ex}
The functions $f\colon \PosInt \to \PosInt$ and $g \colon \PosInt \to
\PosInt$ given by
\begin{align*}
f(n) & = 2n \text{ and}\\
g(n) & = 2n - 1
\end{align*}
enumerate the even positive integers and the odd positive integers,
respectively. However, neither function is an enumeration of
$\PosInt$, since neither is !!{surjective}.
\end{ex}

\begin{prob}
Define an enumeration of the positive squares $1$, $4$, $9$, $16$, \dots
\end{prob}

\begin{ex}
The function $f(n) = (-1)^{n} \lceil \frac{(n-1)}{2}\rceil$ (where
$\lceil x \rceil$ denotes the \emph{ceiling} function, which rounds
$x$ up to the nearest integer) enumerates the set of
integers~$\Int$. Notice how $f$ generates the values of $\Int$ by
``hopping'' back and forth between positive and negative integers:
\[
\begin{array}{c c c c c c c c}
f(1) & f(2) & f(3) & f(4) & f(5) & f(6) & f(7) & \dots \\ \\
- \lceil \tfrac{0}{2} \rceil & \lceil \tfrac{1}{2}\rceil & - \lceil \tfrac{2}{2} \rceil & \lceil \tfrac{3}{2} \rceil & - \lceil \tfrac{4}{2} \rceil  & \lceil \tfrac{5}{2}
\rceil & - \lceil \tfrac{6}{2} \rceil & \dots \\ \\
0 & 1 & -1 & 2 & -2 & 3 & \dots
\end{array}
\]
You can also think of $f$ as defined by cases as follows:
\[
f(n) = \begin{cases}
  0 & \text{if $n = 1$}\\
  n/2 & \text{if $n$ is even}\\
  -(n-1)/2 & \text{if $n$ is odd and $>1$}
  \end{cases}
\]
\end{ex}

\begin{prob}
  Show that if $A$ and $B$ are !!{enumerable}, so is $A \cup B$. To do
  this, suppose there are !!{surjective} functions $f\colon \PosInt \to
  A$ and $g\colon \PosInt \to B$, and define !!a{surjective}
  function~$h\colon \PosInt \to A \cup B$ and prove that it is
  !!{surjective}. Also consider the cases where $A$ or~$B = \emptyset$.
\end{prob}
  
\begin{prob}
  Show that if $B \subseteq A$ and $A$ is !!{enumerable}, so is~$B$. To
  do this, suppose there is !!a{surjective} function $f\colon \PosInt \to
  A$. Define !!a{surjective} function~$g\colon \PosInt \to B$ and prove
  that it is !!{surjective}. What happens if $B = \emptyset$?
\end{prob}
    
\begin{prob}
  Show by induction on $n$ that if $A_1$, $A_2$, \dots, $A_n$ are all
  !!{enumerable}, so is $A_1 \cup \dots \cup A_n$. You may assume the
  fact that if two sets $A$ and~$B$ are !!{enumerable}, so is~$A \cup
  B$. 
\end{prob}

Although it is perhaps more natural when listing the !!{element}s of a
set to start counting from the $1$st !!{element}, mathematicians like
to use the natural numbers~$\Nat$ for counting things. They
talk about the $0$th, $1$st, $2$nd, and so on, !!{element}s of a list.
Correspondingly, we can define an enumeration as !!a{surjective}
function from $\Nat$ to~$A$. Of course, the two definitions are
equivalent.

\begin{prop}\ollabel{prop:enum-shift}
  There is !!a{surjection} $f\colon \PosInt \to A$ iff there is
  !!a{surjection} $g\colon \Nat \to A$.
\end{prop}

\begin{proof}
  Given !!a{surjection} $f\colon \PosInt \to A$, we can define $g(n) =
  f(n+1)$ for all $n \in \Nat$. It is easy to see that $g\colon \Nat
  \to A$ is !!{surjective}. Conversely, given !!a{surjection} $g\colon
  \Nat \to A$, define $f(n) = g(n-1)$.
\end{proof}

This gives us the following result:

\begin{cor}\ollabel{cor:enum-nat}
A set $A$ is !!{enumerable} iff it is empty or there is
!!a{surjective} function $f\colon \Nat \to A$.
\end{cor}

We discussed above that a list of !!{element}s of a set~$A$ can be
turned into a list without repetitions. This is also true for
enumerations, but a bit harder to formulate and prove rigorously. Any
function $f\colon \PosInt \to A$ must be defined for all $n \in
\PosInt$. If there are only finitely many !!{element}s in~$A$ then we
clearly cannot have a function defined on the infinitely many
!!{element}s of~$\PosInt$ that takes as values all the !!{element}s
of~$A$ but never takes the same value twice. In that case, i.e., in
the case where the list without repetitions is finite, we must choose
a different domain for~$f$, one with only finitely many !!{element}s.
Not having repetitions means that $f$ must be !!{injective}. Since it
is also !!{surjective}, we are looking for !!a{bijection} between some
finite set $\{1, \dots, n\}$ or $\PosInt$ and~$A$.

\begin{prop}\ollabel{prop:enum-bij}
If $f\colon \PosInt \to A$ is !!{surjective} (i.e., an enumeration
of~$A$), there is !!a{bijection} $g\colon Z \to A$ where $Z$ is
either~$\PosInt$ or $\{1, \dots, n\}$ for some~$n \in \PosInt$.
\end{prop}

\begin{proof}
  We define the function $g$ recursively: Let $g(1) = f(1)$. If $g(i)$
  has already been defined, let $g(i+1)$ be the first value of $f(1)$,
  $f(2)$, \dots{} not already among $g(1)$, \dots, $g(i)$, if there is
  one. If $A$ has just $n$ !!{element}s, then $g(1)$, \dots, $g(n)$ are all
  defined, and so we have defined a function $g\colon \{1, \dots, n\}
  \to A$. If $A$ has infinitely many !!{element}s, then for any $i$
  there must be !!a{element} of~$A$ in the enumeration $f(1)$, $f(2)$,
  \dots, which is not already among $g(1)$, \dots, $g(i)$. In this
  case we have defined a function $g\colon \PosInt \to A$.
  
  The function $g$ is !!{surjective}, since any element of~$A$ is
  among $f(1)$, $f(2)$, \dots{} (since $f$ is !!{surjective}) and so
  will eventually be a value of~$g(i)$ for some~$i$. It is also
  !!{injective}, since if there were $j < i$ such that $g(j) = g(i)$,
  then $g(i)$ would already be among $g(1)$, \dots, $g(i-1)$, contrary
  to how we defined~$g$.
\end{proof}

\begin{cor}\ollabel{cor:enum-nat-bij}
A set $A$ is !!{enumerable} iff it is empty or there is !!a{bijection}
$f\colon N \to A$ where either $N = \Nat$ or $N = \{0, \dots, n\}$ for
some $n \in \Nat$.
\end{cor}

\begin{proof}
$A$ is !!{enumerable} iff $A$ is empty or there is !!a{surjective}
$f\colon \PosInt \to A$. By \olref{prop:enum-bij}, the latter holds
iff there is !!a{bijective} function~$f\colon Z \to A$ where $Z =
\PosInt$ or $Z = \{1, \dots, n\}$ for some $n \in \PosInt$. By the
same argument as in the proof of \olref{prop:enum-shift}, that in turn
is the case iff there is !!a{bijection} $g\colon N \to A$ where either
$N = \Nat$ or $N = \{0, \dots, n-1\}$.
\end{proof}

\begin{prob}
  According to \olref[sfr][siz][enm]{defn:enumerable}, a set $A$ is
  enumerable iff $A = \emptyset$ or there is !!a{surjective} $f\colon
  \PosInt \to A$.  It is also possible to define ``!!{enumerable} set''
  precisely by: a set is enumerable iff there is !!a{injective}
  function $g\colon A \to \PosInt$.  Show that the definitions are
  equivalent, i.e., show that there is !!a{injective} function
  $g\colon A \to \PosInt$ iff either $A = \emptyset$ or there is
  !!a{surjective} $f\colon \PosInt \to A$.
\end{prob}





\olfileid{sfr}{siz}{zigzag}

\olsection{Cantor's Zig-Zag Method}

\begin{explain}
We've already considered some ``easy'' enumerations. Now we will
consider something a bit harder. Consider the set of pairs of natural
numbers, which we defined in
\olref[set][pai]{sec} thus:
\[
\Nat \times \Nat = \Setabs{\tuple{n,m}}{n,m \in \Nat}
\]
We can organize these ordered pairs into an \emph{array}, like so:
\[
\begin{array}{ c | c | c | c | c | c}
& \mathbf 0 & \mathbf 1 & \mathbf 2 & \mathbf 3 & \dots \\
\hline
\mathbf 0 & \tuple{0,0} & \tuple{0,1} & \tuple{0,2} & \tuple{0,3} & \dots \\
\hline
\mathbf 1 & \tuple{1,0} & \tuple{1,1} & \tuple{1,2} & \tuple{1,3} & \dots \\
\hline
\mathbf 2 & \tuple{2,0} & \tuple{2,1} & \tuple{2,2} & \tuple{2,3} & \dots \\
\hline
\mathbf 3 & \tuple{3,0} & \tuple{3,1} & \tuple{3,2} & \tuple{3,3} & \dots \\
\hline
\vdots & \vdots & \vdots & \vdots & \vdots & \ddots\\
\end{array}
\]
Clearly, every ordered pair in $\Nat \times \Nat$ will appear
exactly once in the array. In particular, $\tuple{n,m}$ will appear in
the $n$th row and $m$th column. But how do we organize the elements of
such an array into a ``one-dimensional'' list? The pattern in the array below
demonstrates one way to do this (although of course there are many other options):
\[
\begin{array}{ c | c | c | c | c | c | c}
& \mathbf 0 & \mathbf 1 & \mathbf 2 & \mathbf 3 & \mathbf 4 &\dots \\
\hline
\mathbf 0 & 0  & 1& 3 & 6& 10 &\ldots \\
\hline
\mathbf 1 &2 & 4& 7 & 11 & \dots &\ldots \\
\hline
\mathbf 2 & 5 & 8 & 12 & \ldots & \dots&\ldots \\
\hline
\mathbf 3 & 9 & 13 & \ldots & \ldots & \dots & \ldots \\
\hline
\mathbf 4 & 14 & \ldots & \ldots & \ldots & \dots & \ldots \\
\hline
\vdots & \vdots & \vdots & \vdots & \vdots&\ldots & \ddots\\
\end{array}
\]\noindent
This pattern is called \emph{Cantor's zig-zag method}. It  enumerates
$\Nat \times \Nat$ as follows:
\[
\tuple{0,0}, \tuple{0,1}, \tuple{1,0}, \tuple{0,2}, \tuple{1,1},
\tuple{2,0}, \tuple{0,3}, \tuple{1,2}, \tuple{2,1}, \tuple{3,0}, \dots
\]
And this establishes the following:
\end{explain}

\begin{prop}\ollabel{natsquaredenumerable}
$\Nat \times \Nat$ is !!{enumerable}.
\end{prop}

\begin{proof}
Let $f \colon \Nat \to \Nat\times\Nat$ take each $k \in \Nat$ to the
tuple $\tuple{n,m} \in \Nat \times \Nat$ such that $k$ is the value of
the $n$th row and $m$th column in Cantor's zig-zag array. 
\end{proof}

\begin{explain}
This technique also generalises rather nicely. For example, we can use
it to enumerate the set of ordered triples of natural numbers, i.e.:
\[
\Nat \times \Nat \times \Nat = \Setabs{\tuple{n,m,k}}{n,m,k \in \Nat}
\]
We think of $\Nat \times \Nat \times \Nat$ as the Cartesian
product of $\Nat \times \Nat$ with $\Nat$, that is,
\[
\Nat^3 = (\Nat \times \Nat) \times \Nat =
\Setabs{\tuple{\tuple{n,m},k}}{n, m, k
  \in \Nat }
\]
and thus we can enumerate $\Nat^3$ with an array by labelling one
axis with the enumeration of $\Nat$, and the other axis with the
enumeration of $\Nat^2$:
\[
\begin{array}{ c | c | c | c | c | c}
& \mathbf 0 & \mathbf 1 & \mathbf 2 & \mathbf 3 & \dots \\
\hline
\mathbf{\tuple{0,0}} & \tuple{0,0,0} & \tuple{0,0,1} & \tuple{0,0,2} & \tuple{0,0,3} & \dots \\
\hline
\mathbf{\tuple{0,1}} & \tuple{0,1,0} & \tuple{0,1,1} & \tuple{0,1,2} & \tuple{0,1,3} & \dots \\
\hline
\mathbf{\tuple{1,0}} & \tuple{1,0,0} & \tuple{1,0,1} & \tuple{1,0,2} & \tuple{1,0,3} & \dots \\
\hline
\mathbf{\tuple{0,2}} & \tuple{0,2,0} & \tuple{0,2,1} & \tuple{0,2,2} & \tuple{0,2,3} & \dots\\
\hline
\vdots & \vdots & \vdots & \vdots & \vdots & \ddots \\
\end{array}
\]
Thus, by using a method like Cantor's zig-zag method, we may similarly
obtain an enumeration of~$\Nat^3$. And we can keep going, obtaining
enumerations of $\Nat^n$ for any natural number $n$. So, we have:
\end{explain}

\begin{prop}
$\Nat^n$ is !!{enumerable}, for every $n \in \Nat$.
\end{prop}

\begin{prob}\label{sfr:siz:zigzag:prob:posint-n}
Show that $(\PosInt)^n$ is !!{enumerable}, for every $n \in \Nat$.
\end{prob}

\begin{prob}\label{sfr:siz:zigzag:prob:posint-star}
  Show that $(\PosInt)^*$ is !!{enumerable}. You may assume \cref{sfr:siz:zigzag:prob:posint-n}.
\end{prob}





\olfileid{sfr}{siz}{pai}
\olsection{Pairing Functions and Codes}

\begin{explain}
Cantor's zig-zag method makes the enumerability of $\Nat^n$ visually
evident. But let us focus on our array depicting $\Nat^2$. Following
the zig-zag line in the array and counting the places, we can check
that $\tuple{1,2}$ is associated with the number~$7$. However, it would
be nice if we could compute this more directly. That is, it would be
nice to have to hand the \emph{inverse} of the zig-zag enumeration,
$g\colon \Nat^2 \to \Nat$, such that
\[
g(\tuple{0,0}) = 0, \;
g(\tuple{0,1}) = 1, \;
g(\tuple{1,0}) = 2, \; \dots,  
g(\tuple{1,2}) = 7, \; \dots
\]
This would enable us to calculate exactly where $\tuple{n, m}$ will occur
in our enumeration. 

In fact, we can define $g$ directly by making two observations. First:
if the $n$th row and $m$th column contains value~$v$, then the
$(n+1)$st row and $(m-1)$st column contains value $v + 1$. Second: the
first row of our enumeration consists of the triangular numbers,
starting with $0$, $1$, $3$, $6$, etc. The $k$th triangular number is
the sum of the natural numbers $< k$, which can be computed as
$k(k+1)/2$. Putting these two observations together, consider this
function:
\[
  g(n,m) = \frac{(n+m+1)(n+m)}{2} + n
\]
We often just write $g(n, m)$ rather that $g(\tuple{n, m})$, since it
is easier on the eyes. This tells you first to determine the
$(n+m)^\text{th}$ triangle number, and then add $n$ to it. And
it populates the array in exactly the way we would like. So in
particular, the pair $\tuple{1, 2}$ is sent to $\frac{4 \times 3}{2} +
1 = 7$. 

This function $g$ is the \emph{inverse} of an enumeration of a set of
pairs. Such functions are called \emph{pairing functions}.
\end{explain}

\begin{defn}[Pairing function] 
  A function $f\colon A \times B \to \Nat$ is an arithmetical
  \emph{pairing function} if $f$ is injective. We also say that $f$
  \emph{encodes} $A \times B$, and that $f(x,y)$ is the
  \emph{code} for $\tuple{x,y}$.
\end{defn}

\begin{explain}
We can use pairing functions to encode, e.g., pairs of natural numbers;
or, in other words, we can represent each \emph{pair} of elements
using a \emph{single} number. Using the inverse of the pairing
function, we can \emph{decode} the number, i.e., find out which
pair it represents.
\end{explain}

\begin{prob}
Give an enumeration of the set of all non-negative rational numbers. 
\end{prob}

\begin{prob}
Show that $\Rat$ is !!{enumerable}. Recall that any rational number 
can be written as a fraction $z/m$ with $z \in \Int$, $m \in \Nat^+$.
\end{prob}

\begin{prob}
Define an enumeration of $\Bin^*$.
\end{prob}

\begin{prob}
Recall from your introductory logic course that each possible truth
table expresses a truth function. In other words, the truth functions
are all functions from $\Bin^k \to \Bin$ for some~$k$. Prove that the
set of all truth functions is enumerable.
\end{prob}

\begin{prob}
Show that the set of all finite subsets of an arbitrary infinite
!!{enumerable} set is !!{enumerable}.
\end{prob}

\begin{prob}
A subset of $\Nat$ is said to be \emph{cofinite} iff it is the
complement of a finite set $\Nat$; that is, $A \subseteq \Nat$ is
cofinite iff $\Nat\setminus A$ is finite. Let $I$ be the set whose
!!{element}s are exactly the finite and cofinite subsets of $\Nat$.
Show that $I$ is !!{enumerable}.
\end{prob}

\begin{prob}
Show that the !!{enumerable} union of !!{enumerable} sets is
!!{enumerable}. That is, whenever $A_1$, $A_2$, \dots{} are sets, and
each $A_i$ is !!{enumerable}, then the union $\bigcup_{i=1}^\infty
A_i$ of all of them is also !!{enumerable}. [NB: this is hard!]
\end{prob}

\begin{prob}
Let $f \colon A \times B \to \Nat$ be an arbitrary pairing function.
Show that the inverse of $f$ is an enumeration of $A \times B$.
\end{prob}

\begin{prob}
Specify a function that encodes $\Nat^3$.
\end{prob}





\olfileid{sfr}{siz}{pai-alt}

\olsection{An Alternative Pairing Function}

\begin{explain}
There are other enumerations of $\Nat^2$ that make it easier to
figure out what their inverses are. Here is one. Instead of
visualizing the enumeration in an array, start with the list of
positive integers associated with (initially) empty spaces. Imagine
filling these spaces successively with pairs $\tuple{n,m}$ as follows.
Starting with the pairs that have~$0$ in  the first place (i.e., pairs
$\tuple{0,m}$), put the first (i.e., $\tuple{0,0}$) in the first empty
place, then skip an empty space, put the second (i.e., $\tuple{0,2}$)
in the next empty place, skip one again, and so forth. The
(incomplete) beginning of our enumeration now looks like this
\[\small
\begin{array}{@{}c c c c c c c c c c c@{}}
\mathbf 1 & \mathbf 2 & \mathbf 3 & \mathbf 4 & \mathbf 5 & \mathbf 6 & \mathbf 7 & \mathbf 8 & \mathbf 9 & \mathbf{10} & \dots \\ \\
\tuple{0,0} &  & \tuple{0,1} &  & \tuple{0,2} &  & \tuple{0,3} & & \tuple{0,4} &  & \dots \\
\end{array}
\]
Repeat this with pairs $\tuple{1,m}$ for the place that still remain
empty, again skipping every other empty place:
\[\small
\begin{array}{@{}c c c c c c c c c c c@{}}
\mathbf 1 & \mathbf 2 & \mathbf 3 & \mathbf 4 & \mathbf 5 & \mathbf 6 & \mathbf 7 & \mathbf 8 & \mathbf 9 & \mathbf{10} & \dots \\ \\
\tuple{0,0} & \tuple{1,0} & \tuple{0,1} &  & \tuple{0,2} & \tuple{1,1} & 
\tuple{0,3} & & \tuple{0,4} &  \tuple{1,2} & \dots \\
\end{array}
\]
Enter pairs $\tuple{2,m}$, $\tuple{2,m}$, etc., in the same way. Our
completed enumeration thus starts like this:
\[\small
\begin{array}{@{}cc c c c c c c c c c@{}}
\mathbf 1 & \mathbf 2 & \mathbf 3 & \mathbf 4 & \mathbf 5 & \mathbf 6 & \mathbf 7 & \mathbf 8 & \mathbf 9 & \mathbf{10} & \dots \\ \\
\tuple{0,0} & \tuple{1,0} & \tuple{0,1} & \tuple{2,0}  & \tuple{0,2} & 
\tuple{1,1} & \tuple{0,3} & \tuple{3,0}  & \tuple{0,4} &  \tuple{1,2} & \dots \\
\end{array}
\]
If we number the cells in the array above according to this
enumeration, we will not find a neat zig-zag line, but this
arrangement:
\[
\begin{array}{ c | c | c | c | c | c | c | c }
& \mathbf 0 & \mathbf 1 & \mathbf 2 & \mathbf 3 & \mathbf 4 & \mathbf 5 & \dots \\
\hline
\mathbf 0 & 1 & 3 & 5 & 7 & 9 & 11 & \dots \\
\hline
\mathbf 1 & 2 & 6 & 10 & 14 & 18 & \dots & \dots \\
\hline
\mathbf 2 & 4 & 12 & 20 & 28 & \dots & \dots & \dots \\
\hline
\mathbf 3 & 8 & 24 & 40 & \dots & \dots & \dots & \dots \\
\hline
\mathbf 4 & 16 & 48 & \dots & \dots & \dots & \dots & \dots \\
\hline
\mathbf 5 & 32 & \dots & \dots & \dots & \dots & \dots & \dots \\
\hline
\vdots & \vdots & \vdots & \vdots & \vdots & \vdots & \vdots & \ddots\\
\end{array}
\]
We can see that the pairs in row~$0$ are in the odd numbered places of
our enumeration, i.e., pair $\tuple{0,m}$ is in place $2m+1$; pairs in
the second row, $\tuple{1,m}$, are in places whose number is the
double of an odd number, specifically,  $2 \cdot (2m+1)$; pairs in the
third row, $\tuple{2,m}$, are in places whose number is four times an
odd number, $4 \cdot (2m+1)$; and so on. The factors of $(2m+1)$ for
each row, $1$, $2$, $4$, $8$, \dots, are exactly the powers of~$2$:
$1= 2^0$, $2 = 2^1$, $4 = 2^2$, $8 = 2^3$, \dots\@ In fact, the
relevant exponent is always the first member of the pair in
question. Thus, for pair $\tuple{n,m}$ the factor is $2^n$.  This
gives us the general formula: $2^n \cdot (2m+1)$. However, this is a
mapping of pairs to \emph{positive} integers, i.e., $\tuple{0,0}$ has
position~$1$. If we want to begin at position~$0$ we must subtract~$1$
from the result. This gives us:
\end{explain}

\begin{ex}
The function $h\colon \Nat^2 \to \Nat$ given by
\[
h(n,m) = 2^n (2m+1) - 1
\]
is a pairing function for the set of pairs of natural numbers~$\Nat^2$.
\end{ex}

\begin{explain}
Accordingly, in our second enumeration of $\Nat^2$, the pair
$\tuple{0,0}$ has code $h(0,0) = 2^0(2\cdot 0+1) - 1 = 0$;
$\tuple{1,2}$ has code $2^{1} \cdot (2 \cdot 2 + 1) - 1 = 2
\cdot 5 - 1 = 9$; $\tuple{2,6}$ has code $2^{2} \cdot (2
\cdot 6 + 1) - 1 = 51$.
\end{explain}

Sometimes it is enough to encode pairs of natural numbers~$\Nat^2$
without requiring that the encoding is surjective. Such encodings have
inverses that are only partial functions. 

\begin{ex}
The function $j\colon \Nat^2 \to \Nat^+$ given by
\[
j(n,m) = 2^n3^m
\]
is !!a{injective} function $\Nat^2 \to \Nat$.
\end{ex}





\olfileid{sfr}{siz}{nen}

\olsection{\printtoken{S}{nonenumerable} Sets}

\begin{editorial}
  This section proves the non-enumerability of $\Bin^\omega$ and
  $\Pow{\PosInt}$ using the definition in \olref[enm]{sec}. It is
  designed to be a little more elementary and a little more detailed
  than the version in \olref[enm-alt]{sec}
\end{editorial}

Some sets, such as the set $\PosInt$ of positive integers, are
infinite. So far we've seen examples of infinite sets which were all
!!{enumerable}. However, there are also infinite sets which do not
have this property. Such sets are called \emph{!!{nonenumerable}}.

First of all, it is perhaps already surprising that there are
!!{nonenumerable} sets.  For any !!{enumerable} set~$A$ there is
!!a{surjective} function $f \colon \PosInt \to A$.  If a set is
!!{nonenumerable} there is no such function.  That is, no function
mapping the infinitely many !!{element}s of~$\PosInt$ to~$A$ can
exhaust all of~$A$.  So there are ``more'' !!{element}s of~$A$ than
the infinitely many positive integers.

How would one prove that a set is !!{nonenumerable}? You have to show
that no such surjective function can exist. Equivalently, you have to
show that the elements of~$A$ cannot be enumerated in a one way
infinite list.  The best way to do this is to show that every list of
!!{element}s of~$A$ must leave at least one element out; or that no
function $f\colon \PosInt \to A$ can be !!{surjective}.  We can do this
using Cantor's \emph{diagonal method}.  Given a list of !!{element}s
of~$A$, say, $x_1$, $x_2$, \dots, we construct another element of~$A$
which, by its construction, cannot possibly be on that list.

Our first example is the set~$\Bin^\omega$ of all infinite, non-gappy
sequences of $0$'s and $1$'s.

\begin{thm}
\ollabel{thm:nonenum-bin-omega}
$\Bin^\omega$~is !!{nonenumerable}.
\end{thm}

\begin{proof}
Suppose, by way of contradiction, that $\Bin^\omega$ is
!!{enumerable}, i.e., suppose that there is a list $s_{1}$, $s_{2}$,
$s_{3}$, $s_{4}$, \dots{} of all !!{element}s of~$\Bin^\omega$.  Each
of these $s_i$ is itself an infinite sequence of $0$'s and~$1$'s.
Let's call the $j$-th element of the $i$-th sequence in this list
$s_i(j)$. Then the $i$-th sequence~$s_i$ is
\[
s_i(1), s_i(2), s_i(3), \dots
\]

We may arrange this list, and the elements of each sequence $s_i$ in
it, in an array:
\[
\begin{array}{c|c|c|c|c|c}
& 1 & 2 & 3 & 4 & \dots \\\hline
1 & \mathbf{s_{1}(1)} & s_{1}(2) & s_{1}(3) & s_1(4) & \dots \\\hline
2 & s_{2}(1)& \mathbf{s_{2}(2)} & s_2(3) & s_2(4) & \dots \\\hline
3 & s_{3}(1)& s_{3}(2) & \mathbf{s_3(3)} & s_3(4) & \dots \\\hline
4 & s_{4}(1)& s_{4}(2) & s_4(3) & \mathbf{s_4(4)} & \dots \\\hline
\vdots & \vdots & \vdots & \vdots & \vdots & \mathbf{\ddots}
\end{array}
\]
The labels down the side give the number of the sequence in the list
$s_1$, $s_2$, \dots; the numbers across the top label the !!{element}s
of the individual sequences. For instance, $s_{1}(1)$ is a name for
whatever number, a $0$ or a~$1$, is the first !!{element} in the
sequence $s_{1}$, and so on.

Now we construct an infinite sequence, $\overline{s}$, of $0$'s and
$1$'s which cannot possibly be on this list.  The definition of
$\overline{s}$ will depend on the list $s_1$, $s_2$, \dots.  Any
infinite list of infinite sequences of $0$'s and $1$'s gives rise to
an infinite sequence~$\overline{s}$ which is guaranteed to not appear
on the list.

To define $\overline{s}$, we specify what all its !!{element}s are,
i.e., we specify $\overline{s}(n)$ for all $n \in \PosInt$.  We do this
by reading down the diagonal of the array above (hence the name
``diagonal method'') and then changing every $1$ to a $0$ and every
$0$ to a~$1$. More abstractly, we define $\overline{s}(n)$ to be $0$
or $1$ according to whether the $n$-th !!{element} of the diagonal,
$s_n(n)$, is $1$ or $0$.
\[
\overline{s}(n) =
\begin{cases}
1 & \text{if $s_{n}(n) = 0$}\\
0 & \text{if $s_{n}(n) = 1$}.
\end{cases}
\]
If you like formulas better than definitions by cases, you could also
define $\overline{s}(n) = 1 - s_n(n)$.

Clearly $\overline{s}$ is an infinite sequence of $0$'s and
$1$'s, since it is just the mirror sequence to the sequence of $0$'s
and $1$'s that appear on the diagonal of our array.  So $\overline{s}$
is !!a{element} of~$\Bin^\omega$.  But it cannot be on the list $s_1$,
$s_2$, \dots{} Why not?

It can't be the first sequence in the list, $s_1$, because it differs from
$s_1$ in the first !!{element}.  Whatever $s_1(1)$ is, we defined
$\overline{s}(1)$ to be the opposite.  It can't be the second
sequence in the list, because $\overline{s}$ differs from $s_2$ in the second
element: if $s_2(2)$ is $0$, $\overline{s}(2)$ is $1$, and vice
versa. And so on.

More precisely: if $\overline{s}$ were on the list, there would be
some $k$ so that $\overline{s} = s_{k}$.  Two sequences are identical
iff they agree at every place, i.e., for any~$n$, $\overline{s}(n) =
s_{k}(n)$.  So in particular, taking $n = k$ as a special case,
$\overline{s}(k) = s_{k}(k)$ would have to hold. $s_k(k)$ is either
$0$ or~$1$. If it is $0$ then $\overline{s}(k)$ must be~$1$---that's
how we defined $\overline{s}$. But if $s_k(k) = 1$ then, again because
of the way we defined $\overline{s}$, $\overline{s}(k) = 0$. In either
case $\overline{s}(k) \neq s_{k}(k)$.

We started by assuming that there is a list of !!{element}s of
$\Bin^\omega$, $s_1$, $s_2$, \dots{} From this list we constructed a
sequence~$\overline{s}$ which we proved cannot be on the list.  But it
definitely is a sequence of $0$'s and $1$'s if all the $s_i$ are
sequences of $0$'s and $1$'s, i.e., $\overline{s} \in
\Bin^\omega$. This shows in particular that there can be no list of
\emph{all} !!{element}s of~$\Bin^\omega$, since for any such list we
could also construct a sequence~$\overline{s}$ guaranteed to not be on
the list, so the assumption that there is a list of all sequences
in~$\Bin^\omega$ leads to a contradiction.
\end{proof}

\begin{explain}
This proof method is called ``diagonalization'' because it uses the
diagonal of the array to define~$\overline{s}$. Diagonalization need
not involve the presence of an array: we can show that sets are not
!!{enumerable} by using a similar idea even when no array and no
actual diagonal is involved.
\end{explain}

\begin{thm}
\ollabel{thm:nonenum-pownat}
$\Pow{\PosInt}$ is not !!{enumerable}.
\end{thm}

\begin{proof}
We proceed in the same way, by showing that for every list of subsets
of~$\PosInt$ there is a subset of $\PosInt$ which cannot be on the list.
Suppose the following is a given list of subsets of~$\PosInt$:
\[
Z_{1}, Z_{2}, Z_{3}, \dots
\]
We now define a set $\overline{Z}$ such that for any $n \in \PosInt$,
$n \in \overline{Z}$ iff $n \notin Z_{n}$:
\[
\overline{Z} = \Setabs{n \in \PosInt}{n \notin Z_n}
\]
$\overline{Z}$ is clearly a set of positive integers, since by
assumption each~$Z_n$ is, and thus $\overline{Z} \in
\Pow{\PosInt}$. But $\overline{Z}$ cannot be on the list.  To show
this, we'll establish that for each $k \in \PosInt$, $\overline{Z} \neq
Z_k$. 

So let $k \in \PosInt$ be arbitrary. We've defined $\overline{Z}$ so
that for any $n \in \PosInt$, $n \in \overline{Z}$ iff $n \notin Z_n$.
In particular, taking $n=k$, $k \in \overline{Z}$ iff $k \notin Z_k$.
But this shows that $\overline{Z} \neq Z_k$, since $k$ is !!a{element}
of one but not the other, and so $\overline{Z}$ and $Z_k$ have
different !!{element}s. Since $k$ was arbitrary, $\overline{Z}$ is not
on the list $Z_1$, $Z_2$, \dots
\end{proof}

\begin{explain}
The preceding proof did not mention a diagonal, but you can think of
it as involving a diagonal if you picture it this way: Imagine the
sets $Z_1$, $Z_2$, \dots, written in an array, where each
!!{element}~$j \in Z_i$ is listed in the~$j$-th column. Say the first
four sets on that list are $\{1,2,3,\dots\}$, $\{2, 4, 6, \dots\}$,
$\{1,2,5\}$, and $\{3,4,5,\dots\}$. Then the array would begin with
\[
\begin{array}{r@{}rrrrrrr}
  Z_1 = \{ & \mathbf{1}, & 2, & 3, & 4, & 5, & 6, & \dots\}\\
  Z_2 = \{ &  & \mathbf{2}, &  & 4, &  & 6, & \dots\}\\
  Z_3 = \{ & 1, & 2, &  &  & 5\phantom{,} &  & \}\\
  Z_4 = \{ &  &  & 3, & \mathbf{4}, & 5, & 6, & \dots\}\\
  \vdots & & & & & \ddots
\end{array}
\]
Then $\overline{Z}$ is the set obtained by going down the diagonal,
leaving out any numbers that appear along the diagonal and include
those $j$ where the array has a gap in the $j$-th row/column. In the
above case, we would leave out $1$ and $2$, include~$3$, leave
out~$4$, etc.
\end{explain}

\begin{prob}
Show that $\Pow{\Nat}$ is !!{nonenumerable} by a diagonal argument.
\end{prob}

\begin{prob}\label{sfr:siz:nen:prob:f-posint}
Show that the set of functions $f \colon \PosInt \to \PosInt$ is
!!{nonenumerable} by an explicit diagonal argument. That is, show that
if $f_1$, $f_2$, \dots, is a list of functions and each $f_i\colon
\PosInt \to \PosInt$, then there is some $\overline{f}\colon \PosInt \to
\PosInt$ not on this list.
\end{prob}





\olfileid{sfr}{siz}{red}

\olsection{Reduction}

\begin{editorial}
  This section proves non-enumerability by reduction, matching the
  results in \olref[nen]{sec}. An alternative, slightly more condensed
  version matching the results in \olref[nen-alt]{sec} is provided in
  \olref[red-alt]{sec}.
\end{editorial}

We showed $\Pow{\PosInt}$ to be !!{nonenumerable} by a diagonalization
argument. We already had a proof that $\Bin^\omega$, the set of all
infinite sequences of $0$s and $1$s, is !!{nonenumerable}.  Here's
another way we can prove that $\Pow{\PosInt}$ is !!{nonenumerable}:
Show that \emph{if $\Pow{\PosInt}$ is !!{enumerable} then $\Bin^\omega$
  is also !!{enumerable}}.  Since we know $\Bin^\omega$ is not
!!{enumerable}, $\Pow{\PosInt}$ can't be either.  This is called
\emph{reducing} one problem to another---in this case, we reduce the
problem of enumerating $\Bin^\omega$ to the problem of enumerating
$\Pow{\PosInt}$.  A solution to the latter---an enumeration of
$\Pow{\PosInt}$---would yield a solution to the former---an enumeration
of $\Bin^\omega$.

How do we reduce the problem of enumerating a set~$B$ to that of
enumerating a set~$A$?  We provide a way of turning an enumeration
of~$A$ into an enumeration of~$B$.  The easiest way to do that is to
define !!a{surjective} function $f\colon A \to B$.  If $x_1$, $x_2$,
\dots{} enumerates~$A$, then $f(x_1)$, $f(x_2)$, \dots{} would
enumerate~$B$.  In our case, we are looking for a surjective
function $f\colon \Pow{\PosInt} \to \Bin^\omega$.

\begin{prob}
Show that if there is an !!{injective} function $g\colon B \to A$, and
$B$~is !!{nonenumerable}, then so is~$A$. Do this by showing how you
can use~$g$ to turn an enumeration of~$A$ into one of~$B$.
\end{prob}

\begin{proof}[Proof of {\olref[nen]{thm:nonenum-pownat}} by reduction]
Suppose that $\Pow{\PosInt}$ were !!{enumerable}, and thus that
there is an enumeration of it, $Z_{1}$, $Z_{2}$, $Z_{3}$, \dots

Define the function $f \colon \Pow{\PosInt} \to \Bin^\omega$ by letting
$f(Z)$ be the sequence $s_{k}$ such that $s_{k}(n) = 1$ iff $n \in Z$,
and $s_k(n) = 0$ otherwise.  This clearly defines a function, since
whenever $Z \subseteq \PosInt$, any $n \in \PosInt$ either is
!!a{element} of $Z$ or isn't.  For instance, the set $2\PosInt = \{2,
4, 6, \dots\}$ of positive even numbers gets mapped to the sequence
$010101\dots$, the empty set gets mapped to $0000\dots$ and the set
$\PosInt$ itself to $1111\dots$.

It also is !!{surjective}: Every sequence of $0$s and $1$s corresponds
to some set of positive integers, namely the one which has as its
members those integers corresponding to the places where the sequence
has~$1$s. More precisely, suppose $s \in \Bin^\omega$.  Define $Z
\subseteq \PosInt$ by:
\[
Z = \Setabs{n \in \PosInt}{s(n) = 1}
\]
Then $f(Z) = s$, as can be verified by consulting the definition
of~$f$.

Now consider the list
\[
f(Z_1), f(Z_2), f(Z_3), \dots
\]
Since $f$ is !!{surjective}, every member of $\Bin^\omega$ must
appear as a value of~$f$ for some argument, and so must appear on the
list. This list must therefore enumerate all of~$\Bin^\omega$.

So if $\Pow{\PosInt}$ were !!{enumerable}, $\Bin^\omega$ would be
!!{enumerable}.  But $\Bin^\omega$ is !!{nonenumerable}
(\olref[nen]{thm:nonenum-bin-omega}). Hence $\Pow{\PosInt}$ is
!!{nonenumerable}.
\end{proof}

\begin{explain}
It is easy to be confused about the direction the reduction goes in.
For instance, !!a{surjective} function $g \colon \Bin^\omega \to B$
does \emph{not} establish that $B$ is !!{nonenumerable}.  (Consider $g
\colon \Bin^\omega \to \Bin$ defined by $g(s) = s(1)$, the function
that maps a sequence of $0$'s and $1$'s to its first !!{element}.  It
is !!{surjective}, because some sequences start with $0$ and some start
with $1$. But $\Bin$ is finite.)  Note also that the function~$f$ must
be !!{surjective}, or otherwise the argument does not go through:
$f(x_1)$, $f(x_2)$, \dots{} would then not be guaranteed to include
all the !!{element}s of~$B$. For instance, 
\[
h(n) = \underbrace{000\dots0}_{\text{$n$ $0$'s}}
\]
defines a function $h\colon \PosInt \to
\Bin^\omega$, but $\PosInt$ is !!{enumerable}.
\end{explain}

\begin{prob}
Show that the set of all \emph{sets of} pairs of positive integers is
!!{nonenumerable} by a reduction argument.
\end{prob}

\begin{prob}\label{sfr:siz:red:prob:nat-nat}
  Show that the set~$X$ of all functions $f\colon \Nat \to \Nat$ is
  !!{nonenumerable} by a reduction argument (Hint: give a surjective
  function from $X$ to~$\Bin^\omega$.)
\end{prob}

\begin{prob}
Show that $\Nat^\omega$, the set of infinite sequences of
natural numbers, is !!{nonenumerable} by a reduction argument.
\end{prob}

\begin{prob}
Let $P$ be the set of functions from the set of positive
integers to the set $\{0\}$, and let $Q$ be the set of \emph{partial}
functions from the set of positive integers to the set $\{0\}$. Show
that $P$~is !!{enumerable} and $Q$~is not. (Hint: reduce the problem
of enumerating $\Bin^\omega$ to enumerating~$Q$).
\end{prob}

\begin{prob}
Let $S$ be the set of all !!{surjective} functions from the set of
positive integers to the set \{0,1\}, i.e., $S$ consists of all
!!{surjective}~$f\colon \PosInt \to \Bin$.  Show that $S$ is
!!{nonenumerable}.
\end{prob}

\begin{prob}
Show that the set~$\Real$ of all real numbers is !!{nonenumerable}.
\end{prob}





\olfileid{sfr}{siz}{equ}

\olsection{Equinumerosity}

We have an intuitive notion of ``size'' of sets, which works fine for
finite sets. But what about infinite sets? If we want to come up with
a formal way of comparing the sizes of two sets of \emph{any} size, it
is a good idea to start by defining when sets are the same size. Here
is Frege:
\begin{quote}
  If a waiter wants to be sure that he has laid exactly as many knives
  as plates on the table, he does not need to count either of them, if
  he simply lays a knife to the right of each plate, so that every
  knife on the table lies to the right of some plate. The plates and
  knives are thus uniquely correlated to each other, and indeed
  through that same spatial relationship. \citep[\S70]{Frege1884}
\end{quote}
The insight of this passage can be brought out through a formal
definition:

\begin{defn}\ollabel{comparisondef}
  $A$ is \emph{equinumerous} with $B$, written $\cardeq{A}{B}$, iff
  there is !!a{bijection} $f \colon A \to B$. 
\end{defn}

\begin{prop}\ollabel{equinumerosityisequi}
Equinumerosity is an equivalence relation.
\end{prop}

\begin{proof} 
We must show that equinumerosity is reflexive, symmetric, and
transitive. Let $A, B$, and $C$ be sets.

\emph{Reflexivity.} The identity map $\Id{A} \colon A \to A$, where
$\Id{A} (x) = x$ for all $x \in A$, is !!a{bijection}. So
$\cardeq{A}{A}$.

\emph{Symmetry.} Suppose $\cardeq{A}{B}$, i.e., there is
!!a{bijection} $f\colon A \to B$. Since $f$ is !!{bijective}, its
inverse $f^{-1}$ exists and is also !!{bijective}. Hence,
$f^{-1}\colon B \to A$ is !!a{bijection}, so $\cardeq{B}{A}$.

\emph{Transitivity.} Suppose that $\cardeq{A}{B}$ and $\cardeq{B}{C}$,
i.e., there are !!{bijection}s $f\colon A \to B$ and $g\colon B \to
C$. Then the composition $\comp{f}{g}\colon A \to C$ is !!{bijective},
so that $\cardeq{A}{C}$.
\end{proof}

\begin{prop}
If $\cardeq{A}{B}$, then $A$ is !!{enumerable} if
and only if $B$ is.
\end{prop}

\begin{editorial}
The following proof uses \olref[enm]{defn:enumerable} if
\olref[enm]{sec} is included and \olref[enm-alt]{defn:enumerable}
otherwise.
\end{editorial}

\begin{proof}
Suppose $\cardeq{A}{B}$, so there is some !!{bijection} $f \colon A
\to B$, and suppose that $A$ is !!{enumerable}.

  Then either $A = \emptyset$ or there is !!a{surjective} function
  $g\colon \PosInt \to A$. If $A = \emptyset$, then $B = \emptyset$
  also (otherwise there would be !!a{element}~$y \in B$ but no $x \in
  A$ with $g(x) = y$). If, on the other hand, $g\colon \PosInt \to A$
  is !!{surjective}, then $\comp{g}{f} \colon \PosInt \to B$ is
  !!{surjective}. To see this, let $y \in B$. Since $f$ is
  !!{surjective}, there is an $x \in A$ such that $f(x) = y$. Since
  $g$ is !!{surjective}, there is an $n \in \PosInt$ such that $g(n) =
  x$. Hence,
\[
(\comp{g}{f})(n) = f(g(n)) = f(x) = y
\]
and thus $\comp{g}{f}$ is !!{surjective}. We have that $\comp{g}{f}$
is an enumeration of~$B$, and so $B$~is !!{enumerable}.

If $B$ is !!{enumerable}, we obtain that $A$ is !!{enumerable} by
repeating the argument with the !!{bijection} $f^{-1}\colon B \to A$
instead of~$f$. 
\end{proof}

\begin{prob}
Show that if $\cardeq{A}{C}$ and $\cardeq{B}{D}$, and $A \cap B =
C \cap D = \emptyset$, then $\cardeq{A \cup B}{C \cup D}$.
\end{prob}

\begin{prob}
Show that if $A$ is infinite and !!{enumerable}, then
$\cardeq{A}{\Nat}$.
\end{prob}





\olfileid{sfr}{siz}{car}

\olsection{Sets of Different Sizes, and Cantor's Theorem}

\begin{explain}
We have offered a precise statement of the idea that two sets have the
same size. We can also offer a precise statement of the idea that one
set is smaller than another. Our definition of ``is smaller than (or
equinumerous)'' will require, instead of !!a{bijection} between the
sets, !!a{injection} from the first set to the second. If such a
function exists, the size of the first set is less than or equal to
the size of the second. Intuitively, !!a{injection} from one set to
another guarantees that the range of the function has at least as many
!!{element}s as the domain, since no two !!{element}s of the domain
map to the same !!{element} of the range.
\end{explain}

\begin{defn}
$A$ is \emph{no larger than}~$B$, written $\cardle{A}{B}$, iff there
is !!a{injection} $f \colon A \to B$.
\end{defn}

It is clear that this is a reflexive and transitive relation, but that
it is not symmetric (this is left as an exercise). We can also
introduce a notion, which states that one set is (strictly) smaller
than another. 

\begin{defn}
$A$ is \emph{smaller than}~$B$, written $\cardless{A}{B}$, iff there
is !!a{injection}~$f\colon A \to B$ but no !!{bijection}~$g\colon A
\to B$, i.e., $\cardle{A}{B}$ and $\cardneq{A}{B}$.
\end{defn}

It is clear that this relation is irreflexive
and transitive. (This is left as an exercise.) Using this notation, we
can say that a set $A$ is !!{enumerable} iff $\cardle{A}{\Nat}$, and
that $A$ is !!{nonenumerable} iff $\cardless{\Nat}{A}$. This allows us
to restate

\olref[sfr][siz][nen-alt]{thm:nonenum-pownat}
as the observation that
$\cardless{\Nat}{\Pow{\Nat}}$. In fact,
\citet{Cantor1892} proved that this last point is \emph{perfectly
general}:

\begin{thm}[Cantor]\ollabel{thm:cantor}
$\cardless{A}{\Pow{A}}$, for any set $A$.
\end{thm}

\begin{proof}
The map $f(x) = \{x\}$ is !!a{injection} $f \colon A \to \Pow{A}$,
since if $x \neq y$, then also $\{x\} \neq \{y\}$ by extensionality,
and so $f(x) \neq f(y)$. So we have that $\cardle{A}{\Pow{A}}$. 

\begin{editorial}
We present the slow proof if \olref[nen]{sec} is
present, otherwise a faster proof matching \olref[nen-alt]{sec}.
\end{editorial}


  We will now show that there cannot be !!a{surjective} function~$g\colon A \to
  \Pow{A}$, let alone !!a{bijective} one, and hence that
  $\cardneq{A}{\Pow{A}}$. For suppose that $g\colon A \to \Pow{A}$.
  Since $g$ is total, every $x \in A$ is mapped to a subset $g(x)
  \subseteq A$. We can show that $g$ cannot be surjective. To do this, we
  define a subset~$\overline{A} \subseteq A$ which by definition cannot be in the
  range of~$g$. Let
  \[
  \overline{A} = \Setabs{x \in A}{x \notin g(x)}.
  \]
  Since $g(x)$ is defined for all $x \in A$, $\overline{A}$ is clearly
  a well-defined subset of~$A$.  But, it cannot be in the range
  of~$g$. Let $x \in A$ be arbitrary, we will show that $\overline{A} \neq
  g(x)$.  If $x \in g(x)$, then it does not satisfy $x \notin g(x)$,
  and so by the definition of~$\overline{A}$, we have $x \notin
  \overline{A}$.  If $x \in \overline{A}$, it must satisfy the
  defining property of~$\overline{A}$, i.e., $x \in A$ and $x \notin
  g(x)$. Since $x$ was arbitrary, this shows that for each $x \in
  \overline{A}$, $x \in g(x)$ iff $x \notin \overline{A}$, and so
  $g(x) \neq \overline{A}$.  In other words, $\overline{A}$ cannot be
  in the range of~$g$, contradicting the assumption that~$g$ is
  surjective.{}
\end{proof}

\begin{explain}
It's instructive to
  compare the proof of \olref{thm:cantor} to that of
  \olref[nen]{thm:nonenum-pownat}. There we showed that for any list
  $Z_1$, $Z_2$, \dots, of subsets of~$\PosInt$ one can construct a
  set~$\overline{Z}$ of numbers guaranteed not to be on the list. It
  was guaranteed not to be on the list because, for every $n \in
  \PosInt$, $n \in Z_n$ iff $n \notin \overline{Z}$. This way, there
  is always some number that is !!a{element} of one of $Z_n$ or
  $\overline{Z}$ but not the other. We follow the same idea here,
  except the indices~$n$ are now !!{element}s of~$A$ instead
  of~$\PosInt$. The set $\overline{A}$ is defined so that it is
  different from~$g(x)$ for each $x \in A$, because $x \in g(x)$ iff
  $x \notin \overline{A}$. Again, there is always !!a{element} of~$A$
  which is !!a{element} of one of $g(x)$ and $\overline{A}$ but not
  the other. And just as $\overline{Z}$ therefore cannot be on the
  list $Z_1$, $Z_2$, \dots, $\overline{A}$ cannot be in the range
  of~$g$.

It's instructive to
  compare the proof of \olref{thm:cantor} to that of
  \olref[nen-alt]{thm:nonenum-pownat}. There we showed that for any
  list $N_0$, $N_1$, $N_2$, \dots, of subsets of~$\Nat$ we can construct a
  set~$D$ of numbers guaranteed not to be on the list. It was
  guaranteed not to be on the list because $n \in N_n$ iff $n \notin
  D$, for every $n \in \Nat$. We follow the same idea here, except the
  indices~$n$ are now !!{element}s of~$A$ rather than of~$\Nat$. The
  set $D$ is defined so that it is different from~$g(x)$ for each $x
  \in A$, because $x \in g(x)$ iff $x \notin D$.

The proof is also worth comparing with the proof of Russell's Paradox,
\olref[sfr][set][rus]{thm:russells-paradox}. Indeed, Cantor's Theorem was
the inspiration for Russell's own paradox.
\end{explain}

\begin{prob}
  Show that there cannot be !!a{injection} $g\colon \Pow{A} \to
  A$, for any set~$A$. Hint: Suppose $g\colon \Pow{A} \to A$ is
  !!{injective}. Consider $D = \Setabs{g(B)}{B \subseteq A \text{ and
  } g(B) \notin B}$. Let $x = g(D)$. Use the fact that $g$ is
  !!{injective} to derive a contradiction. 
\end{prob}





\olfileid{sfr}{siz}{sb}

\olsection{The Notion of Size, and Schr\"oder-Bernstein}

\begin{explain}
Here is an intuitive thought: if $A$ is no larger than $B$ and $B$ is
no larger than $A$, then $A$ and $B$ are equinumerous. To be honest,
if this thought were \emph{wrong}, then we could scarcely justify the
thought that our defined notion of equinumerosity has anything to do
with comparisons of ``sizes'' between sets!{} Fortunately, though,
the intuitive thought is correct. This is justified by the
Schr\"oder-Bernstein Theorem.
\end{explain}

\begin{thm}[Schr\"oder-Bernstein]
	\ollabel{thm:schroder-bernstein}
	If $\cardle{A}{B}$ and $\cardle{B}{A}$,
	then $\cardeq{A}{B}$.
\end{thm}

\begin{explain}
In other words, if there is !!a{injection} from $A$ to~$B$, and
!!a{injection} from $B$ to~$A$, then there is !!a{bijection} from $A$
to~$B$. 

This result, however, is really rather \emph{difficult} to prove.
Indeed, although Cantor stated the result, others proved
it.\footnote{For more on the history, see e.g.,
\citet[pp.~165--6]{Potter2004}.}

For now, you can (and must)
take it on trust. 

Fortunately, Schr\"oder-Bernstein is \emph{correct}, and it
vindicates our thinking of the relations we defined, i.e.,
$\cardeq{A}{B}$ and $\cardle{A}{B}$, as having something to do with
``size''. Moreover, Schr\"oder-Bernstein is very \emph{useful}. It
can be difficult to think of !!a{bijection} between two equinumerous
sets. The Schr\"oder-Bernstein Theorem allows us to break the comparison
down into cases so we only have to think of !!a{injection} from the
first to the second, and vice-versa.
\end{explain}


\begin{editorial}
The following \olref[sfr][siz][enm-alt]{sec},
\olref[sfr][siz][nen-alt]{sec}, \olref[sfr][siz][red-alt]{sec} are
alternative versions of \olref[sfr][siz][enm]{sec},
\olref[sfr][siz][nen]{sec}, \olref[sfr][siz][red]{sec} due to Tim
Button for use in his Open Set Theory text. They are slightly more
advanced and use a difference definition of enumerability more
suitable in a set theory context (i.e., bijection with $\Nat$ or an
initial segment, rather than being listable or being the range of a
surjective function from $\PosInt$).
\end{editorial}



\olfileid{sfr}{siz}{enm-alt}

\olsection{Enumerations and \usetoken{S}{enumerable} Sets}

\begin{editorial}
  This section defines enumerations as bijections with (initial
  segments) of $\Nat$, the way it's done in set theory. So it
  conflicts slightly with the definitions in \olref[enm]{sec}, and
  repeats all the examples there. It is also a bit more terse than
  that section.
\end{editorial}

We can specify finite set is by simply enumerating its
!!{element}s. We do this when we define a set like so:
\[
  A = \{a_1, a_2, \ldots, a_n\}.
\]
Assuming that the !!{element}s $a_1$, \dots, $a_n$ are all distinct,
this gives us !!a{bijection} between $A$ and the first $n$ natural
numbers $0$, \dots, $n-1$. Conversely, since every finite set has only
finitely many !!{element}s, every finite set can be put into such a
correspondence. In other words, if $A$ is finite, there is
!!a{bijection} between $A$ and $\{0, \dots, n-1\}$, where $n$ is the
number of !!{element}s of~$A$.

If we allow for certain kinds of infinite sets, then we will also
allow some infinite sets to be enumerated. We can make this precise by
saying that an infinite set is enumerated by !!a{bijection} between it
and all of~$\Nat$.

\begin{defn}[Enumeration, set-theoretic] 
An \emph{enumeration} of a set $A$ is !!a{bijection} whose range is
$A$ and whose domain is either an initial set of natural numbers $\{0,
1, \ldots, n\}$ {or} the entire set of natural numbers~$\Nat$. 
\end{defn}

\begin{explain}
There is an intuitive underpinning to this use of the word
\emph{enumeration}. For to say that we have enumerated a set $A$ is to
say that there is !!a{bijection} $f$ which allows us to count out the
elements of the set $A$. The $0$th element is $f(0)$, the 1st is
$f(1)$, \ldots the $n$th is $f(n)$\ldots.\footnote{Yes, we count
from $0$. Of course we could also start with~$1$. This would
make no big difference. We would just have to replace~$\Nat$
by~$\PosInt$.} The rationale for this may be made even clearer by
adding the following:
\end{explain}

\begin{defn}
  \ollabel{defn:enumerable}
  A set~$A$ is !!{enumerable} iff either $A = \emptyset$ or there is
  an enumeration of~$A$. We say that $A$ is !!{nonenumerable} iff $A$
  is not !!{enumerable}.
\end{defn}

\begin{explain}
So a set is !!{enumerable} iff it is empty or you can use an
enumeration to count out its !!{element}s.
\end{explain}

\begin{ex}
A function enumerating the natural numbers is simply the identity
function $\Id{\Nat} \colon \Nat \to \Nat$ given by $\Id{\Nat}(n) = n$. A
function enumerating the \emph{positive} natural numbers, $\Nat^+ =
\Nat \setminus \{0\}$, is the function $g(n) = n + 1$, i.e., the
successor function.
\end{ex}

\begin{prob}
Show that a set $A$ is !!{enumerable} iff either $A = \emptyset$ or
there is !!a{surjection} $f\colon \Nat \to A$. Show that $A$ is
!!{enumerable} iff there is !!a{injection} $g\colon A \to \Nat$. 
\end{prob}

\begin{ex}
The functions $f\colon \Nat \to \Nat$ and $g \colon \Nat \to \Nat$
given by
\begin{align*}
  f(n) & = 2n \text{ and}\\
  g(n) & = 2n+1
\end{align*}
respectively enumerate the even natural numbers and the odd natural
numbers. But neither is !!{surjective}, so neither is an enumeration
of $\Nat$.
\end{ex}

\begin{prob}
Define an enumeration of the square numbers $1$, $4$, $9$, $16$, \dots
\end{prob}

\begin{ex}
Let $\lceil x \rceil$ be the \emph{ceiling} function, which rounds $x$
up to the nearest integer. Then the function $f \colon \Nat \to \Int$
given by:
\[
  f(n) = (-1)^{n} \left\lceil\tfrac{n}{2}\right\rceil
\]
enumerates the set of
integers~$\Int$ as follows:
\[
\begin{array}{c c c c c c c c}
f(0) & f(1) & f(2) & f(3) & f(4) & f(5) & f(6) & \dots \\ \\
\big\lceil \tfrac{0}{2} \big\rceil & -\big\lceil \tfrac{1}{2}\big\rceil &  \big\lceil \tfrac{2}{2} \big\rceil & -\big\lceil \tfrac{3}{2} \big\rceil & \big\lceil \tfrac{4}{2} \big\rceil  & -\big\lceil \tfrac{5}{2}\big\rceil & \big\lceil \tfrac{6}{2} \big\rceil & \dots \\ \\
0 & -1 & 1 & -2 & 2 & -3 & 3& \dots
\end{array}
\]
Notice how $f$ generates the values of $\Int$ by ``hopping'' back and
forth between positive and negative integers. You can also think of
$f$ as defined by cases as follows:
\[
f(n) = \begin{cases}
  \frac{n}{2} & \text{if $n$ is even}\\
  -\frac{n+1}{2} & \text{if $n$ is odd}
  \end{cases}
\]
\end{ex}

\begin{prob}
Show that if $A$ and $B$ are !!{enumerable}, so is $A \cup B$.
\end{prob}

\begin{prob}
Show by induction on $n$ that if $A_1$, $A_2$, \dots, $A_n$ are all
!!{enumerable}, so is $A_1 \cup \dots \cup A_n$.
\end{prob}




\olfileid{sfr}{siz}{nen-alt}

\olsection{\printtoken{S}{nonenumerable} Sets}

\begin{editorial}
  This section proves the non-enumerability of $\Bin^\omega$ and
  $\Pow{\Nat}$ using the definitions in \olref[enm-alt]{sec}, i.e.,
  requiring a bijection with~$\Nat$ instead of a surjection from
  $\PosInt$.
\end{editorial}

\begin{explain}
The set $\Nat$ of natural numbers is infinite. It is also trivially
!!{enumerable}. But the remarkable fact is that there are
\emph{!!{nonenumerable}} sets, i.e., sets which are not !!{enumerable}
(see \olref[sfr][siz][enm-alt]{defn:enumerable}).

This might be surprising. After all, to say that $A$ is
!!{nonenumerable} is to say that there is \emph{no} !!{bijection} $f
\colon \Nat \to A$; that is, no function mapping the infinitely many
!!{element}s of~$\Nat$ to~$A$ exhausts all of~$A$.  So if $A$ is
!!{nonenumerable}, there are ``more'' !!{element}s of~$A$ than there
are natural numbers.

To prove that a set is !!{nonenumerable}, you have to show that no
appropriate !!{bijection} can exist. The best way to do this is to
show that every attempt to enumerate !!{element}s of~$A$ must leave at
least one !!{element} out; this shows that no function $f\colon \Nat
\to A$ is !!{surjective}. And a general strategy for establishing this
is to use Cantor's \emph{diagonal method}. Given a list of
!!{element}s of $A$, say, $x_1$, $x_2$, \dots, we construct another
!!{element} of~$A$ which, by its construction, cannot possibly be on
that list.

But all of this is best understood by example. So, our first example
is the set~$\Bin^\omega$ of all infinite strings of $0$'s and $1$'s.
(The `$\Bin$' stands for binary, and we can just think of it as the
two-element set
$\{0,1\}$.)\end{explain}

\begin{thm}
\ollabel{thm:nonenum-bin-omega}
$\Bin^\omega$~is !!{nonenumerable}.
\end{thm}

\begin{proof}
Consider any enumeration of a subset of $\Bin^\omega$. So we have some
list $s_{0}$, $s_{1}$, $s_{2}$, \dots{} where every $s_n$ is an
infinite string of $0$'s and~$1$'s. Let $s_n(m)$ be the $n$th digit of
the $m$th string in this list. So we can now think of our list as an
array, where $s_n(m)$ is placed at the $n$th row and $m$th column:
\[
\begin{array}{c|c|c|c|c|c}
& 0 & 1 & 2 & 3 & \dots \\\hline
0 & \mathbf{s_{0}(0)} & s_{0}(1) & s_{0}(2) & s_0(3) & \dots \\\hline
1 & s_{1}(0)& \mathbf{s_{1}(1)} & s_1(2) & s_1(3) & \dots \\\hline
2 & s_{2}(0)& s_{2}(1) & \mathbf{s_2(2)} & s_2(3) & \dots \\\hline
3 & s_{3}(0)& s_{3}(1) & s_3(2) & \mathbf{s_3(3)} & \dots \\\hline
\vdots & \vdots & \vdots & \vdots & \vdots & \mathbf{\ddots}
\end{array}
\]
We will now construct an infinite string, $d$, of $0$'s and $1$'s
which is not on this list.  We will do this by specifying each of its
entries, i.e., we specify $d(n)$ for all $n \in \Nat$.  Intuitively,
we do this by reading down the diagonal of the array above (hence the
name ``diagonal method'') and then changing every $1$ to a $0$ and
every $1$ to a~$0$. More abstractly, we define $d(n)$ to be $0$ or $1$
according to whether the $n$-th !!{element} of the diagonal, $s_n(n)$,
is $1$ or $0$, that is:
\[
d(n) =
\begin{cases}
1 & \text{if $s_{n}(n) = 0$}\\
0 & \text{if $s_{n}(n) = 1$}
\end{cases}
\]
Clearly $d \in \Bin^\omega$, since it is an infinite string of $0$'s
and $1$'s. But we have constructed $d$ so that $d(n) \neq s_n(n)$ for
any $n \in \Nat$. That is, $d$ differs from $s_n$ in its $n$th entry.
So $d \neq s_n$ for any $n\in \Nat$. So $d$ cannot be on the list
$s_0$, $s_1$, $s_2$,
\dots

We have shown, given an arbitrary enumeration of some subset of
$\Bin^\omega$, that it will omit some !!{element} of $\Bin^\omega$. So
there is no enumeration of the set $\Bin^\omega$, i.e., $\Bin^\omega$
is !!{nonenumerable}.
\end{proof}

\begin{explain}
This proof method is called ``diagonalization'' because it uses the
diagonal of the array to define~$d$. However, diagonalization need
not involve the presence of an array. Indeed, we can show that some set is 
!!{nonenumerable} by using a similar idea, even when no array and no
actual diagonal is involved. The following result illustrates how.
\end{explain}

\begin{thm}
\ollabel{thm:nonenum-pownat}
$\Pow{\Nat}$ is not !!{enumerable}.
\end{thm}

\begin{proof}
We proceed in the same way, by showing that every list of subsets
of~$\Nat$ omits some subset of $\Nat$. So, suppose that we have some
list $N_0, N_1, N_2, \ldots$ of subsets of $\Nat$. We define a set $D$
as follows: $n \in D$ iff $n \notin N_{n}$:
\[
D = \Setabs{n \in \Nat}{n \notin N_n}
\]
Clearly $D\subseteq \Nat$. But $D$ cannot be on the list. After all,
by construction $n \in D$ iff $n\notin N_n$, so that $D \neq N_n$ for
any $n \in \Nat$. 
\end{proof}

\begin{explain}
The preceding proof did not mention a diagonal. Still, you can think
of it as involving a diagonal if you picture it this way: Imagine the
sets $N_0$, $N_1$, \dots, written in an array, where we write $N_n$ on
the $n$th row by writing $m$ in the $m$th column iff if $m \in N_n$.
For example, say the first four sets on that list are
$\{0,1,2,\dots\}$, $\{1, 3, 5, \dots\}$, $\{0,1,4\}$, and
$\{2,3,4,\dots\}$; then our array would begin with
\[
\begin{array}{r@{}rrrrrrr}
  N_0 = \{ & \mathbf{0}, & 1, & 2, & & & & \dots\}\\
  N_1 = \{ &  & \mathbf{1}, &  & 3, &  & 5, & \dots\}\\
  N_2 = \{ & 0, & 1, &  &  & 4\phantom{,} &  & \}\\
  N_3 = \{ &  &  & 2, & \mathbf{3}, & 4, & & \dots\}\\
  &\vdots & & & & & & \ddots\phantom{\}}
  \end{array}
\]
Then $D$ is the set obtained by going down the diagonal, placing $n
\in D$ iff $n$ is \emph{not} on the diagonal. So in the above case, we
would leave out $0$ and $1$, we would include~$2$, we would leave
out~$3$, etc.
\end{explain}

\begin{prob}
Show that the set of all functions $f \colon \Nat \to \Nat$ is
!!{nonenumerable} by an explicit diagonal argument. That is, show that
if $f_1$, $f_2$, \dots, is a list of functions and each $f_i\colon
\Nat \to \Nat$, then there is some $g \colon \Nat \to
\Nat$ not on this list.
\end{prob}




\olfileid{sfr}{siz}{red-alt}

\olsection{Reduction}

\begin{editorial}
  This section proves non-enumerability by reduction, matching the
  results in \olref[nen-alt]{sec}. An alternative, slightly more
  elaborate version matching the results in \olref[nen]{sec} is
  provided in \olref[red]{sec}.
\end{editorial}

We proved that $\Bin^\omega$ is !!{nonenumerable} by a diagonalization
argument. We used a similar diagonalization argument to show that
$\Pow{\Nat}$ is !!{nonenumerable}. But here's another way we can prove
that $\Pow{\Nat}$ is !!{nonenumerable}: show that \emph{if
$\Pow{\Nat}$ is !!{enumerable} then $\Bin^\omega$ is also
!!{enumerable}}.  Since we know $\Bin^\omega$ is !!{nonenumerable}, it
will follow that $\Pow{\Nat}$ is too.  

This is called \emph{reducing} one problem to another. In this case,
we reduce the problem of enumerating $\Bin^\omega$ to the problem of
enumerating $\Pow{\Nat}$.  A solution to the latter---an enumeration
of $\Pow{\Nat}$---would yield a solution to the former---an
enumeration of $\Bin^\omega$.

To reduce the problem of enumerating a set~$B$ to that of enumerating
a set~$A$, we provide a way of turning an enumeration of~$A$ into an
enumeration of~$B$.  The easiest way to do that is to define
!!a{surjection} $f\colon A \to B$.  If $x_1$, $x_2$, \dots{}
enumerates~$A$, then $f(x_1)$, $f(x_2)$, \dots{} would enumerate~$B$.
In our case, we are looking for !!a{surjection} $f\colon \Pow{\Nat}
\to \Bin^\omega$.

\begin{prob}
Show that if there is an !!{injective} function $g\colon B \to A$, and
$B$~is !!{nonenumerable}, then so is~$A$. Do this by showing how you
can use~$g$ to turn an enumeration of~$A$ into one of~$B$.
\end{prob}

\begin{proof}[Proof of {\olref[nen-alt]{thm:nonenum-pownat}} by reduction]
For a reduction, suppose that $\Pow{\Nat}$ is !!{enumerable}, and thus that
there is an enumeration of it, $N_{1}$, $N_{2}$, $N_{3}$, \dots

Define the function $f \colon \Pow{\Nat} \to \Bin^\omega$ by letting
$f(N)$ be the string $s_{k}$ such that $s_{k}(n) = 1$ iff $n \in N$,
and $s_k(n) = 0$ otherwise.  

This clearly defines a function, since whenever $N \subseteq \Nat$,
any $n \in \Nat$ either is !!a{element} of $N$ or isn't.  For
instance, the set $2\Nat = \Setabs{2n}{n \in \Nat} = \{0,2, 4, 6,
\dots\}$ of even naturals gets mapped to the string $1010101\dots$;
$\emptyset$ gets mapped to $0000\dots$; $\Nat$ gets mapped to
$1111\dots$.

It is also !!{surjective}: every string of $0$s and $1$s corresponds
to some set of natural numbers, namely the one which has as its
members those natural numbers corresponding to the places where the string
contains a~$1$s. More precisely, if $s \in \Bin^\omega$, then define $N
\subseteq \Nat$ by:
\[
N = \Setabs{n \in \Nat}{s(n) = 1}
\]
Then $f(N) = s$, as can be verified by consulting the definition
of~$f$. 

Now consider the list
\[
f(N_1), f(N_2), f(N_3), \dots
\]
Since $f$ is !!{surjective}, every member of $\Bin^\omega$ must
appear as a value of~$f$ for some argument, and so must appear on the
list. This list must therefore enumerate all of~$\Bin^\omega$.

So if $\Pow{\Nat}$ were !!{enumerable}, $\Bin^\omega$ would be
!!{enumerable}.  But $\Bin^\omega$ is !!{nonenumerable}
(\olref[nen-alt]{thm:nonenum-bin-omega}). Hence $\Pow{\Nat}$ is
!!{nonenumerable}.
\end{proof}



















\begin{prob}\label{sfr:siz:red:prob:nat-nat}
  Show that the set~$X$ of all functions $f\colon \Nat \to \Nat$ is
  !!{nonenumerable} by a reduction argument (Hint: give a surjective
  function from $X$ to~$\Bin^\omega$.)
\end{prob}

\begin{prob}
Show that the set of all \emph{sets of} pairs of natural numbers,
i.e., $\Pow{\Nat \times \Nat}$, is !!{nonenumerable} by a reduction
argument.
\end{prob}

\begin{prob}
Show that $\Nat^\omega$, the set of infinite sequences of natural
numbers, is !!{nonenumerable} by a reduction argument.
\end{prob}








\begin{prob}
Let $S$ be the set of all !!{surjection}s from $\Nat$ to the set
$\{0,1\}$, i.e., $S$ consists of all !!{surjection}s~$f \colon \Nat
\to \Bin$.  Show that $S$ is !!{nonenumerable}.
\end{prob}

\begin{prob}
Show that the set~$\Real$ of all real numbers is !!{nonenumerable}.
\end{prob}



\OLEndChapterHook

